\chapter{All orders of the directed meridian profile}
\label{ch:directed-all-orders}

The directed criterion gives an expansion of the actual meridian state
to every finite order. Its coefficients satisfy a degree bound that
links the exact quantum recurrence to the preferred geometric
longitude. This extends the geometric-factor conclusion to all the
new directed cases, including the short-power examples permitted by
large degree. The first correction is the cusp shape of the same
filled knot complement.

Throughout this chapter the ordered tree, its signed powers, endpoint
frames, peripheral markings and directed certificate are fixed before
the level tends to infinity. There is at least one edge. The original
strand test must give one knot, with no component closed inside a
tangle. Both level parities are retained. The full proof, including
the translated parameter regions and the uniform remainders, is in
Appendix~\ref{app:directed-all-orders}.

\section{The expansion and its geometric consequences}
\label{dao-ch:statement}

Let the tree satisfy the finite directed inequalities of
Appendix~\ref{app:directed-circle-messages}. This includes the
vertexwise sufficient condition
\begin{equation}\label{dao-ch:sufficient}
 d_v\ge2^{28}\quad\text{or}\quad |n_v|\ge2^{27}
                  \qquad\text{at every vertex}.
\end{equation}
The choice can differ from vertex to vertex. More general directed
certificates, with different subtree budgets and a freely chosen
final vertex, remain included.

Write the actual zero-framed reduced state as
\begin{equation}\label{dao-ch:state}
 q_N=e^{\pi i/N},\qquad M=-q_Nx=q_N^s,\qquad
 K_N=F_N(1),\qquad
 H_N(s)=\frac{\sin\pi s}{N\sin(\pi s/N)}.
\end{equation}
Here $x$ is the original highest-weight parameter. The normalized
full state is $\mathcal Z_N(s)=F_N(q_N^s)/(K_NH_N(s))$.
The one factor $H_N$ is fixed by the reduced cut and modified trace;
it is not an adjustable normalization.

\begin{theorem}[All orders for a directed certificate]
For every fixed compact $K$ in the open upper half-plane and every
finite $L$,
\begin{equation}\label{dao-ch:expansion}
 \log\mathcal Z_N(s)
  =\sum_{r=1}^{L}N^{-r}p_r(s)+O_{K,L}(N^{-L-1}),
       \qquad \deg p_r\le r+1.
\end{equation}
The same polynomials give the expansion on every fixed disk
$|s|\le R<1$. Fixed derivatives are controlled on smaller compacts.
The logarithm is the branch tending to zero. Lower amplitudes may be
organized separately by level parity; the selected leading action
and its top polynomial degrees are the same.

Let $\tau_K$ be the preferred cusp shape of the actual filled knot.
Then
\begin{equation}\label{dao-ch:cusp}
 \frac{F_N(q_N^s)}{K_NH_N(s)}
     =1+\frac{\pi i\tau_Ks^2}{2N}+O_K(N^{-2}).
\end{equation}
For every genuine fixed all-level reduced annihilator, its first
nonzero classical symbol is divisible by the selected complex
geometric meridian--longitude factor.
\end{theorem}

The constants and starting levels in the expansion may depend on
the fixed tree, certificate, compact and requested order. The
numerical directed criterion is unchanged. The volume already
identified by that criterion is the volume of the original filled
knot complement; no auxiliary comparison manifold replaces it here.

\section{Why the actual source has every order}
\label{dao-ch:source}

After lowering the physical trace pairing, every original turn edge
carries its exact paired weight. The full edge kernel is
\begin{equation}\label{dao-ch:kernel}
 \mathscr K_N(t,u;s)=\sum_{k=0}^{N-1}
 e_k(-2\cos t)e_k(-2\cos u)\frac{H_N(s)}{W_k(s)},
\end{equation}
where $e_k$ is the fixed complete Newton polynomial and
\[
 W_k(s)=\prod_{r=1}^{k}
       \left(1-\frac{\sin^2(\pi s/N)}{\sin^2(\pi r/N)}\right).
\]
The original six-type edge scalar is retained outside the kernel.
Each exact subtree output is an even trigonometric polynomial of
degree at most $N-1$ in its remaining angle. Choosing the last
vertex changes the order of finite contractions, not the physical
cut or the framing.

The directed proof normalizes every child by its exact value at
angle zero. It supplies a nonzero conditional first Fourier saddle,
a holomorphic nonzero leading amplitude, and strict exponential
bounds for the remaining charges and all endpoint currents. The
outgoing real and strip estimates are estimates of the exact
polynomial, including its high indices and repeated-node normals.

These facts permit an induction to any prescribed finite order.
At a local saddle, expand the canonical factorials and amplitude,
integrate the finite Gaussian Taylor expansion, and divide by the
corresponding expansion at parent angle zero. The integral of the
absolute local Gaussian is bounded by a fixed multiple of the
absolute leading integral. A uniform relative input remainder
therefore remains a relative remainder of the same order. The
constant can depend on the fixed subtree; no uniform condition
number over all diagrams is asserted.

Outside the local saddle, retain the exact child polynomials and
their original real and strip bounds. A truncated expansion is not
used there. The strict directed margins make the charge tails,
original endpoints and connector terms smaller than every inverse
power of $N$. Restoring all exact child constants recovers the
original joint action and its nonzero amplitude. This proves higher
orders from the finite source, rather than by differentiating its
first-order estimate.

An operator of order $d$ needs $s,s+1,\ldots,s+d$, so the proof must
also cover translated regions. Euler's product expresses $W_k$ as
a subset product of $1-s^2/m^2$ over distinct nonzero integers $m$.
On a fixed nonreal compact both the product and its inverse are
bounded uniformly in $k,N$. The full fixed-power crossing polynomial
has the same polynomial norm bound there. Its common Cartan scalar
and mixed reference factors remain in the exact scalar calculation.
The constant opposite-color channel retains its separate full-normal
limit.

The translated grids move their endpoints by $O(N^{-1})$. The
original endpoint currents and any wrapped grid entries are kept;
the crossing chirp is not declared periodic. The exact charge
coefficient is
\[
 (-1)^{(N-1)(\ell-1)}U_{\ell-1}(\cos2\pi\sigma),
        \qquad \sigma=s,-s,0,
\]
whose first value is one. Its fixed-compact growth in $\ell$ is
absorbed by the existing geometric tail. Thus the same conditional
proof works on every required translated compact, with the same
exponential margins. It asserts no uniform continuation across an
ordinary-color resonance.

\section{The polynomial filtration and the preferred cusp}
\label{dao-ch:filtration}

The degree bound in \eqref{dao-ch:expansion} comes from an exact
insertion. Let $G_N$ be the canonical sine-factorial interpolation,
with $G_N(k)=\prod_{r=1}^{k}2\sin(\pi r/N)$ and its fixed reflection
normalization. On an original wire $x_e=k/N$,
\begin{equation}\label{dao-ch:insertion}
 \frac{H_N(s)}{W_k(s)}
  =\frac{G_N(k)^2}{G_N(k-s)G_N(k+s)}.
\end{equation}
No periodic interpolation factor is allowed. The midpoint offsets
from $N-1$ and the original integer wire offsets remain in its
expansion.

Let $\Phi(X)$ be the original joint first-charge phase, with $X$
collecting the physical wire and circle-angle variables. Put
\begin{equation}\label{dao-ch:phase}
\begin{aligned}
 \mathcal K(x,t)&=2J(x)-J(x+t)-J(x-t),\qquad
 J'(x)=\log(2\sin\pi x),\\
 \Phi_t(X)&=\Phi(X)+\sum_{e\in\mathcal E_T}\mathcal K(x_e,t)
                   +\frac{\pi i a(D)}2t^2,\\
 a(D)&=-w(D)-\sum_vn_v,\qquad S(t)=\Phi_t(X_*(t)).
\end{aligned}
\end{equation}
Here $\mathcal E_T$ is the original physical turn-edge set. The
critical point $X_*(t)$ is the analytic branch through the filled
first-charge point. The Hessian is invertible because each original
wire Hessian and each conditional angular Schur complement is
invertible.

The finite Stirling expansion of $G_N$, to arbitrary requested
order, gives holomorphic functions $C_j$ with
\begin{equation}\label{dao-ch:organized}
 \log\mathcal Z_N(s)\sim
 N\{S(s/N)-S(0)\}
   +\sum_{j\ge0}N^{-j}\{C_j(s/N)-C_j(0)\}.
\end{equation}
This is proved using a local auxiliary family with independent
$t$, then substituting $t=s/N$. It makes no physical asymptotic
claim for $s$ of order $N$. At order $N^{-r}$ the phase contributes
degree at most $r+1$ in $s$, while the lower amplitudes contribute
degree at most $r$. This proves the filtration and identifies its
top diagonal.

Both the phase and the canonically interpolated amplitude are even
in $t$. Normalization removes a constant at order $N^{-1}$ and
evenness removes a possible linear term from $C_0$. Therefore
$p_1(s)=S''(0)s^2/2$.

The geometric identification uses the actual filled branch, not
this derivative alone. The directed certificate places every
restored angle in the common positive charge-path chart. Its
primitive equations remain $U_v+p_vV_v=-2\pi i$ with
$n_v=\epsilon_v-2p_v$. The reconstructed representation is the
complete filled-knot representation, and its nondegenerate
meridian continuation is the geometric germ. The preferred
longitude theorem gives
\begin{equation}\label{dao-ch:longitude}
 L_{\mathrm{geom}}(M)
 =-M^{a(D)}\prod_{e\in\mathcal E_T}
            \frac{M^2-X_e(M)}{1-M^2X_e(M)},
       \qquad X_e=e^{2\pi ix_e}.
\end{equation}
All factors refer to the original diagram and preferred framing.
The negative lift is fixed at the actual complete filled knot.

With $u=\log M$ and
$\mathcal S(u)=\pi i\{S(u/(\pi i))-S(0)\}$,
\eqref{dao-ch:phase} and the exact sine identity give
$-e^{\mathcal S'(u)}=L_{\mathrm{geom}}(e^u)$.
It follows that $S''(0)=\pi i\tau_K$, proving
\eqref{dao-ch:cusp} with the exact framing restored.

\section{The factor forced in a genuine classical symbol}
\label{dao-ch:operator}

Let $P=\sum_{j=0}^{d}a_j(q,M)T^j$, with $TM=qMT$, be a specified
nonzero fixed rational-coefficient operator annihilating the whole
reduced function $F_N(M)$ at every sufficiently large level.
Clear fixed denominators and the common $q-1$ valuation, and denote
its first nonzero classical symbol by $P_0(M,L)$. This is an
all-level identity, not a recurrence fitted to complete values.

The modified dimension is
\[
 d_N(M)=i^{N+3}\frac{M-M^{-1}}{M^N-M^{-N}},
       \qquad \frac{d_N(q_N^s)}{d_N(1)}=H_N(s)^{-1}.
\]
For $f(M)=M-M^{-1}$, the normalized full state obeys the exact gauge
\begin{equation}\label{dao-ch:gauge}
 \widetilde P=\sum_{j=0}^{d}a_j(q,M)(-1)^j
                 \frac{f(M)}{f(q^jM)}T^j.
\end{equation}
Clearing these additional denominators retains the first symbol
$f(M)^dP_0(M,-L)$.

Set $h=\log q$ and regrade by $u=hs$. The degree filtration gives
$h^{-1}\mathcal S(u)+\sum_{j\ge0}h^jA_j(u)$ for the logarithm.
A $j$-fold shift has formal leading ratio $e^{j\mathcal S'(u)}$.
The exact equation and the expansions on all translated disks imply
\begin{equation}\label{dao-ch:symbol}
 f(e^u)^dP_0(e^u,-e^{\mathcal S'(u)})=0.
\end{equation}
The substitution $u\mapsto hs$ is injective on the formal two-variable
series ring. Its domain property permits cancellation of the
nonzero series $f(e^u)^d$ in this proved identity. The meridian
factor has therefore been recorded before cancellation.

Evaluation on the geometric germ defines a nonzero height-one
prime of $\mathbb C[M^{\pm1},L^{\pm1}]$. Its irreducible generator
$A_{\mathrm{geom},\mathbb C}$ divides $P_0$. Over
$\mathbb Q(q,M)$, the conclusion is instead stated using the
contraction of that prime to the rational Laurent ring; its
complex factorization may include conjugate components. No chosen
complex component is asserted to be defined over $\mathbb Q$.

The conclusion determines a selected factor, not its multiplicity,
the full A-polynomial or an equality of operator ideals. It adds no
moving-main theorem for cyclic Morse sources and no fixed-meridian
WKB limit. Those questions require additional source estimates.
